% Bagian: first-order-logic
% Bab: introduction
% Subbagian: semantic-notions

\documentclass[../../../include/open-logic-section]{subfiles}

\begin{document}

\olfileid[id]{fol}{int}{sem}

\olsection{Gagasan Semantis}

Di atas kita menyebutkan bahwa ketika mempertimbangkan apakah
$\Sat{M}{!A}[s]$ berlaku, kita (demi kemudahan) membiarkan $s$ menetapkan
nilai kepada semua !!{variable}s, tetapi hanya nilai yang ditetapkannya kepada
!!{variable}s dalam~$!A$ yang digunakan. Bahkan, hanya nilai
variabel-variabel \emph{bebas} dalam~$!A$ yang berpengaruh. Tentu saja, karena
kita teliti, kita akan membuktikan fakta ini. Karena !!{sentence}s tidak
mempunyai variabel bebas, $s$ sama sekali tidak berpengaruh terhadap apakah
kalimat-kalimat itu terpenuhi dalam !!a{structure}. Jadi, jika $!A$ merupakan
!!a{sentence}, kita dapat mendefinisikan $\Sat{M}{!A}$ sebagai
``$\Sat{M}{!A}[s]$ bagi semua~$s$,'' yang ternyata benar jika dan hanya jika
$\Sat{M}{!A}[s]$ bagi setidaknya satu~$s$. Kita perlu memperkenalkan penugasan
!!{variable} agar memperoleh definisi keterpenuhan yang berfungsi bagi
!!{formula}s, tetapi bagi !!{sentence}s, keterpenuhan tidak bergantung pada
penugasan !!{variable}.

Setelah mempunyai definisi ``$\Sat{M}{!A}$,'' kita mengetahui arti ``kasus''
dan ``benar dalam'' sejauh menyangkut !!{sentence}s logika orde pertama.
Berdasarkan definisi $\Sat{M}{!A}$ bagi !!{sentence}s, kita kemudian dapat
mendefinisikan gagasan semantis dasar berupa validitas, perikutan, dan
keterpenuhan. Suatu kalimat valid, $\Entails !A$, jika setiap !!{structure}
memenuhinya. Kalimat itu diakibatkan oleh suatu himpunan !!{sentence}s,
$\Gamma \Entails !A$, jika setiap !!{structure} yang memenuhi semua
!!{sentence}s dalam~$\Gamma$ juga memenuhi~$!A$. Dan suatu himpunan
!!{sentence}s dapat dipenuhi jika ada suatu !!{structure} yang memenuhi semua
!!{sentence}s di dalamnya sekaligus.

Karena !!{formula}s didefinisikan secara induktif, dan keterpenuhan pada
gilirannya didefinisikan melalui induksi pada struktur !!{formula}s, kita
dapat menggunakan induksi untuk membuktikan sifat-sifat semantik kita dan
menghubungkan gagasan-gagasan semantis yang telah didefinisikan. Kita akan
mengumpulkan dan membuktikan beberapa sifat ini, sebagian karena masing-masing
sifat itu menarik, tetapi terutama karena banyak di antaranya akan berguna
ketika kita melanjutkan dengan menyelidiki hubungan antara semantik dan sistem
!!{derivation}. Untuk melakukannya, kita juga harus mendefinisikan (secara
tepat, yakni melalui induksi) beberapa gagasan dan operasi sintaktis yang
belum kita sebutkan.

\end{document}
